\begin{python0}
from solutions import *; clear()
\end{python0}
\sectionthree{Use of \texttt{NULL} for memory management}

Recall that \texttt{NULL} is a predefined constant for 0 (as an address)
and it's used to denote an invalid address. We can use \texttt{NULL} to
denote the fact that a pointer has not been allocated.

Look at this:

\begin{consolethree}[escapeinside=||]
int * p = NULL;
// *p is not used yet

...

// Now you're ready to use an int on the heap
p = new int;
// use *p

...
// Now you don't need the int on the heap anymore
delete p;
p = NULL;

...
\end{consolethree}
If you always set your pointer \texttt{p} to \texttt{NULL} when memory is not allocated for \texttt{p} to point to, then you can always check your pointer against \texttt{NULL} like this:

\begin{consolethree}[escapeinside=||]
int * p = NULL;

...

if (p == NULL)
{
    p = new int;
}

// use *p ...
delete p;
p = NULL;

...

if (p == NULL)
{
    p = new int;
}

// use *p ...
delete p;
p = NULL;

...
\end{consolethree}
Why use the free store?

You might ask \ldots why bother requesting for an \texttt{int} inside the free store?

\begin{consolethree}[escapeinside=||]
int * p = new int;
// ... now use *p
delete p;
\end{consolethree}
That's so much trouble \ldots why not just create a plain \texttt{int} variable:

\begin{consolethree}[escapeinside=||]
int x;
// ... now use x
\end{consolethree}
That is true!

Actually the notes from the previous section is just to show you how to use the memory heap. In most cases, you use the heap by requesting for a \texttt{\textbf{\underline{very huge}}} chunk of memory from the free store and not for a single \texttt{int}. In other words, dynamic memory management is useful only for large memory usage. As an example, consider the following scenario:

Suppose you have some type \texttt{T} that requires a huge amount of memory. Say a variable of type \texttt{T} is some kind of an image for a game and a variable of type \texttt{T} holds one megabytes of pixel data. Maybe you need 3 such variables:

\begin{center}
\texttt{T bigvar0, bigvar1, bigvar2;}
\end{center}

In this case all these 3 huge variables will consume 3 MB of memory.

But what if you don't always need them at the same time? Say \ldots

\begin{consolethree}[escapeinside=||]
T bigvar0, bigvar1, bigvar2;

// use bigvar0 for 10 min

// use bigvar1 for 10 min

// use bigvar1 and bigvar2 for 10 min

// use bigvar0 for 10 min 
\end{consolethree}
Then you might want to do this instead:

\begin{consolethree}[escapeinside=||]
T *bigvar0, *bigvar1, *bigvar2;

bigvar0 = new T; // 1 MB used
// use bigvar0 for 10 min
delete bigvar0; // 0 MB used

bigvar1 = new T; // 1 MB used
// use bigvar1 for 10 min

bigvar2 = new T; // 2 MB used
// use bigvar1 and bigvar2 for 10 min
delete bigvar1; // 1 MB used
delete bigvar2; // 0 MB used

bigvar0 = new T; // 1 MB used
// use bigvar0 for 10 min
delete bigvar0; // 0 MB used
\end{consolethree}
This means that you actually only need about 2 MB of memory at any point
in time while running the program, not 3 MB.

Get it?

Later I'll talk about requesting not just memory for a value, but a whole array of values. When the array size is huge, this will take up a lot or memory even if each value in the array is only an \texttt{int}. For instance:

\texttt{int x[1000000];}

takes up 4 x 1,000,000 = 4,000,000 bytes. Dynamic memory management is also important when it comes to managing objects (see CISS245, CISS350) since objects usually occupy more memory.

